\hypertarget{oppgave3}{}
\section*{Oppgave 3}
\addcontentsline{toc}{section}{Oppgave 3}

I denne oppgaven genereres firkantpulser og bitsekvenser i Octave.
For å gjøre dette er det laget egne funksjoner for generering av
signalene. Selve implementasjonen av funksjonene er lagt i
vedlegg A, mens denne delen viser hvordan funksjonene brukes og
resultatene de gir.
\begin{verbatim}
pkg load signal; # for å bruke square() funksjonen
\end{verbatim}

% --------------------------------------------------
\subsection*{a)}

For generering av firkantpulsen er funksjonen
\texttt{mySquareV1} laget. Funksjonen bruker signalets fase i
radianer og modulo $2\pi$ for å bestemme hvor i perioden hvert
sample befinner seg. \\[1em]
Det er også laget en alternativ funksjon, \texttt{mySquareV2},
som utfører tilsvarende generering direkte i tidsdomenet.
Implementasjonen av begge funksjonene finnes i vedlegg A.\\[1em]
Det brukes en fundamental frekvens på 10 kHz og 100 samples
per periode:
\begin{verbatim}
f0 = 10e3;
T = 1/f0;
fs = f0*100;

t_start = 0;
dt = 1/fs;
t_end = T*4-dt;

d_c = 50; # %
d_c2 = 20; # %

t = t_start:dt:t_end;
\end{verbatim}
Periodetiden er:
\[
    T = \frac{1}{f_0} = \frac{1}{10kHz} = 0.1\,\text{ms}
\]
Med 100 samples per periode blir samplingsfrekvensen:
\[
    f_s = 100f_0 = 1\,\text{MHz}
\]
Den egenlagde funksjonen brukes først til å generere en puls med
50 \% duty cycle. Deretter genereres en sentrert puls med
20 \% duty cycle:
\begin{verbatim}
[_, out_m] = mySquareV1(t, f0, d_c);
[_, out_m2] = mySquareV1(t, f0, d_c2, [0,1], 0, true);
\end{verbatim}
For å kontrollere resultatet sammenlignes signalene med Octaves
innebygde \texttt{square()}-funksjon:
\begin{verbatim}
out = square(t*2*pi*f0, d_c);

t_1 = T*(d_c2/100)/2;

out2 = square((t+t_1)*2*pi*f0, d_c2);
out2 = (out2 + 1)/2;
\end{verbatim}
Figuren viser at den egenlagde funksjonen gir samme firkantpulser
som \texttt{square()}-funksjonen.
\begin{figure}[H]
    \centering
    \includegraphics[width=1\textwidth]{oppgave3a.png}
    \caption{Sammenligning mellom den egenlagde
    \texttt{mySquareV1}-funksjonen og \texttt{square()} i Octave.}
\end{figure}

% --------------------------------------------------
\subsection*{b)}
For å generere en bitsekvens er funksjonen
\texttt{bitWaveGen} laget. Funksjonen mottar en bitvektor,
bittid, antall samples per bit og et eventuelt shift. Funksjonen bruker også hjelpefunksjonen
\texttt{addZerosOrOnes}. Implementasjonen av begge funksjonene finnes i vedlegg A.
Bitsekvensen som skal genereres er:
\[
[1,\ 0,\ 1,\ 0].
\]
Bittiden settes til 1 ms, med 100 samples per bit og et shift på 50 \%:
\begin{verbatim}
bits = [1,0,1,0];
T = 1e-3; # s
spp = 100; # sample per periode
sh = 50; # %

[t, out] = bitWaveGen(bits, T, spp, sh);
\end{verbatim}
Samplingsintervallet blir:
\[
dt = \frac{T}{100}
   = 10\,\mu\text{s}.
\]
Figuren under viser den resulterende bitsekvensen.
\begin{figure}[H]
    \centering
    \includegraphics[width=0.9\textwidth]{oppgave3b.png}
    \caption{Generert bitsekvens $[1,0,1,0]$ med
    100 samples per bit og 50 \% shift.}
\end{figure}
